\subsection{Construction of continuous sections of étale maps}

THEOREM 1. --- Let X, Y, Z be topological spaces, let \(f: \mathrm{Z} \to \mathrm{X} \times \mathrm{Y}\) be a separated étale morphism, and let \(y_0\) be a point of Y. Let \(s: X \times Y \to Z\) be a section of \(f\) . Assume that the restriction of \(s\) to \(X \times \{y_0\}\) is continuous, as are the restrictions of \(s\) to \(\{x\} \times Y\) for all \(x \in X\). If the space \(Y\) is connected and locally connected (TG, I, p.~84, def. 4), the map \(s\) is continuous.

Lemma. --- Let U be an open subset of X, let V be a connected open subset of Y, and let \((x_{1},y_{1})\in\mathrm{U}\times\mathrm{V}\). Assume that the restriction of s to \(U\times\{y_{1}\}\) is continuous. Let \(\sigma\) be a continuous section of f over \(U\times V\) such that \(\sigma(x_{1},y_{1})=s(x_{1},y_{1})\) . There exists a neighborhood \(U'\) of \(x_{1}\) such that \(s=\sigma\) on \(U'\times V\). In particular, the map s is continuous in a neighborhood of \((x_{1},y_{1})\) .

Since the restriction of \(s\) to \(\mathrm{U} \times \{y_1\}\) is continuous and the map \(f\) is étale, there exists a neighborhood \(\mathrm{U}'\) of \(x_1\) such that \(s(x, y_1) = \sigma(x, y_1)\) for all \(x \in \mathrm{U}'\) (I, p.~34, prop. 11, b)). Let \(x \in \mathrm{U}'\) ; since the restriction of \(s\) to \(\{x\} \times \mathrm{V}\) is continuous and the map \(f\) is étale and separated, it follows from cor. 1, I, p.~34 that \(\sigma(x, y) = s(x, y)\) for all \(y \in \mathrm{V}\) . Thus, \(s\) and \(\sigma\) coincide on \(\mathrm{U}' \times \mathrm{V}\) .

Let us now prove the theorem. Let us first prove that the map s is continuous in a neighborhood of every point of \(X \times \{y_{0}\}\) . Let \(x_{0} \in X\) ; one can choose an open neighborhood U of \(x_{0}\) in X, a connected open neighborhood V of \(y_{0}\) in Y and a continuous section \(\sigma: U \times V \to Z\) of f such that \(\sigma(x_{0}, y_{0}) = s(x_{0}, y_{0})\) . By the lemma above, s is continuous in a neighborhood of \((x_{0}, y_{0})\) .

We now prove that the map \(s\) is continuous at every point of \(\mathrm{X} \times \mathrm{Y}\) . For \(x_0 \in \mathrm{X}\) , denote by \(\mathrm{C}(x_0)\) the set of points \(y \in \mathrm{Y}\) such that \(s\) is continuous in a neighborhood of \((x_0, y)\) . The set \(\mathrm{C}(x_0)\) is open in Y by definition, and contains \(y_0\) by the preceding.

Let us prove that it is closed in Y. Consider a point \(y_{1}\) of Y adherent to \(\mathrm{C}(x_{0})\) , an open neighborhood U of \(x_{0}\) in X, an open neighborhood V of \(y_{1}\) in Y, and a continuous section \(\sigma\colon U\times V\to Z\) of f such that \(\sigma(x_{0},y_{1})=s(x_{0},y_{1})\) . Since the restriction of s to \(\{x_{0}\}\times V\) is continuous and f is étale, there exists an open neighborhood \(V^{\prime}\) of \(y_{1}\) in Y such that \(\sigma(x_{0},y)=s(x_{0},y)\) for all \(y\in V^{\prime}\) (prop. 11 of I, p.~34). Since Y is locally connected, we may assume that \(V^{\prime}\) is connected. Since \(y_{1}\) is adherent to \(\mathrm{C}(x_{0})\), the set \(\mathrm{C}(x_{0})\cap\mathrm{V}^{\prime}\) is nonempty; let \(y_{2}\) be a point of \(\mathrm{C}(x_{0})\cap\mathrm{V}^{\prime}\) . By the lemma applied to \((x_{0},y_{2})\) , there exists a neighborhood \(U^{\prime}\) of \(x_{0}\) contained in U such that \(s=\sigma\) over \(U^{\prime} \times V^{\prime}\) . This proves that \(\mathrm{C}(x_{0})\) is closed, because \(U^{\prime} \times V^{\prime}\) is a neighbourhood of \((x_{0}, y_{1})\) . Since Y is connected, we have \(\mathrm{C}(x_{0}) = \mathrm{Y}\) , and this for every \(x_{0} \in X\) , whence the theorem.

COROLLARY 1. --- Let X, Y, E, B be topological spaces. Let \(p\colon E \to B\) be an étale and separated map, and let \(h\colon X \times Y \to B\) be a continuous map and let \(y_0\) be a point of Y. Let \(g\colon X \times Y \to E\) be a map such that \(p \circ g = h\) . Suppose that the restrictions of \(g\) to \(X \times \{y_0\}\) on the one hand and, for every \(x \in X\) , to \(\{x\} \times Y\) on the other hand, are continuous. If the space Y is connected and locally connected, the map \(g\) is continuous.

Let Z be the fibered product \((\mathrm{X} \times \mathrm{Y}) \times_{\mathrm{B}} \mathrm{E}\) and \(f: Z \to X \times Y\) and \(h': Z \to E\) be the canonical projections. The map f is étale (I, p.~31, prop. 8) and separated (I, p.~27, prop. 4). The map \(s: X \times Y \to Z\) defined by \(s(x, y) = ((x, y), g(x, y))\) is a section of f which satisfies the hypotheses of Theorem 1. It is therefore continuous, as is \(g = h' \circ s\) .

Remark. --- If, in Theorem 1, one does not assume that the space Y is locally connected, the conclusion of the theorem is no longer necessarily valid (I, p.~141, exerc. 7).

THEOREM 2. --- Let B be a topological space and A a subspace of B. Assume one of the following hypotheses:

\begin{enumerate}
\def\labelenumi{(\roman{enumi})}
\item
  A has a fundamental system of paracompact neighborhoods;
\item
  B is paracompact and A is closed;
\item
  B is metrizable;
\item
  A is compact and any two distinct points of A have disjoint neighborhoods in B.
\end{enumerate}

Then, every continuous section over A of an étale map \(p\colon E \to B\) extends to a continuous section of \(p\) over a neighborhood of A.

Consider the following property (PCV) of the pair (B, A):

(PCV) For every covering \((\mathrm{U}_{i})_{i\in\mathrm{I}}\) of A by open subsets of B, there exists a neighborhood V of A and a locally finite family \((\mathrm{F}_{j})_{j\in\mathrm{J}}\) of closed subsets of V covering V, such that each of the \(F_{j}\) is contained in one of the \(U_{i}\) .

Theorem 2 follows from Lemmas 1 and 3 below.

Lemma 1. --- Each of the properties (i) to (iv) of Theorem 2 implies property (PCV).

Let \((\mathrm{U}_{i})_{i\in\mathrm{I}}\) be a covering of A by open subsets of B. Under hypothesis (i), there exists a paracompact neighborhood V of A contained in the union of the \(U_{i}\) , an open covering \((\mathrm{U}_{j}^{\prime})_{j\in\mathrm{J}}\) of the space V, locally finite and finer than the covering \((\mathrm{V}\cap\mathrm{U}_{i})_{i\in\mathrm{I}}\) , and an open covering \((\mathrm{U}_{j}^{\prime\prime})_{j\in\mathrm{J}}\) of the space V such that for every \(j\in J\) , the closure \(F_{j}\) of \(U_{j}^{\prime\prime}\) in V is contained in \(U_{j}^{\prime}\) (TG, IX, p.~49, prop. 4 and p.~48, cor. 1 to th. 3). Hence (PCV) in this case.

Condition (ii) implies condition (i) by Corollary 2 of TG, IX, p.~50. Likewise, condition (iii) implies condition (i): indeed, if B is metrizable, every neighborhood of A is metrizable, hence paracompact (TG, IX, p.~51, th. 4).

Finally suppose condition (iv) is satisfied and let us show that property (PCV) holds. Since A is compact, it suffices to consider the case of a finite covering \((\mathrm{U}_{i})_{0\leqslant i\leqslant n}\) . Let us then construct by induction open sets \(V_{0},\ldots,V_{n}\) of B satisfying the following conditions for \(0\leqslant i\leqslant n\) :

\(\alpha)\mathrm{A}\subset \mathrm{V}_{0}\cup \dots \cup \mathrm{V}_{i}\cup \mathrm{U}_{i + 1}\cup \dots \cup \mathrm{U}_{n};\)

\(\beta)\overline{\mathrm{V}_{i}}\cap \mathrm{A}\subset \mathrm{U}_{i}.\)

Suppose \(V_{0},\ldots,V_{r-1}\) have been constructed satisfying the above conditions for \(0\leqslant i\leqslant r-1\) , and let us construct \(V_{r}\) . The sets \(K=A\cap C U_{r}\) and \(L=A\cap C(V_{0}\cup\cdots\cup V_{r-1}\cup U_{r+1}\cup\cdots\cup U_{n})\) are closed in A, hence compact. By virtue of \((\alpha)\) , they are disjoint. By hypothesis, for every point a of L and every point b of K, there exist disjoint neighborhoods of a and b in B. By Lemma 2 below, there exist open sets \(V_{r}\) and W in B, disjoint and such that \(L\subset V_{r}\) and \(K\subset W\) . From the inclusions \(L\subset V_{r}\) and \(A\subset V_{0}\cup\cdots\cup V_{r-1}\cup U_{r}\cup\cdots\cup U_{n}\) and the definition of L, one deduces \((\alpha)\) for i=r. On the other hand, one has \(\overline{V_{r}}\cap K=\varnothing\) , whence \(\overline{V_{r}}\cap A\subset U_{r}\) .

The set \(M = \bigcup_{0 \leqslant i \leqslant n} (\overline{\mathrm{V}_{i}} \cap \mathbb{C}\mathrm{U}_{i})\) is closed and does not meet A by ( \(\beta\) ). By ( \(\alpha\) ), the set \(V = (\bigcup_{0 \leqslant i \leqslant n} \overline{\mathrm{V}_{i}}) - M\) is a neighborhood of A in B. For \(i = 0, \ldots, n\) , set \(F_{i} = V \cap \overline{V_{i}}\); it is a closed subset of V, contained in \(U_{i}\) . The family \((F_{i})_{0 \leqslant i \leqslant n}\) is a covering of V, hence property (PCV) holds.

Lemma 2. --- Let B be a topological space, and let K and L be quasi-compact subsets of B. Suppose that for every point a of K and every point b of L, there exist disjoint neighborhoods of a and b in B. Then there exist two disjoint open sets U and V in B such that K ⊂ U and L ⊂ V.

Let \(a\) be a point of K. For every \(b \in L\) , let \(U_{a,b}\) and \(V_{a,b}\) be disjoint open neighborhoods of \(a\) and of \(b\) . For \(a\) fixed, the family \((V_{a,b})_{b \in L}\) is an open cover of L. Since L is quasi-compact, there exists a finite family \(T_a \subset L\) such that the union \(V_b\) of the \(V_{a,b}\) for \(b \in T_a\) contains L. The intersection \(U_a\) of the \(U_{a,b}\) for \(b \in T_a\) is an open neighborhood of \(a\) since \(T_a\) is finite; moreover, \(U_a\) and \(V_a\) are disjoint. The \((U_a)_{a \in A}\) form an open covering of K. Since K is quasi-compact, there exists a finite family \(S \subset K\) such that \(U = \bigcup_{a \in S} U_a\) contains K. Then, \(V = \bigcap_{a \in S} V_a\) is an open subset of B containing L. The lemma is proved.

Lemma 3. --- Let B be a topological space and A a subspace of B such that the pair (B,A) satisfies property (PCV) of I, p.~37. Then, every continuous section over A of an étale map \(p\colon E \to B\) extends to a continuous section of \(p\) over a neighborhood of A.

Let \(p \colon \mathrm{E} \to \mathrm{B}\) be an étale map and \(s \colon \mathrm{A} \to \mathrm{E}\) a continuous section of \(p\) over A. For every point \(a\) of A, there exists an open neighborhood \(\mathrm{U}_a\) of \(a\) and a continuous section \(s_a \colon \mathrm{U}_a \to \mathrm{E}\) which coincides with \(s\) on \(\mathrm{U}_a \cap \mathrm{A}\) (I, p.~34, prop. 10 and I, p.~34, prop. 11). The family \((\mathrm{U}_a)_{a \in \mathrm{A}}\) is a covering of A by open sets of B. Let V be a neighborhood of A in B, let \((\mathrm{F}_j)_{j \in \mathrm{J}}\) be a locally finite family of closed subsets of V covering V, and let \(a \colon \mathrm{J} \to \mathrm{A}\) be a map such that \(\mathrm{F}_j\) is contained in \(\mathrm{U}_{a(j)}\) for all \(j \in \mathrm{J}\) (property (PCV)). Let us denote by \(s_j\) the restriction of \(s_{a(j)}\) to \(\mathrm{F}_j\) . Let W be the set of points \(b \in \mathrm{V}\) satisfying \(s_j(b) = s_k(b)\) for every pair \((j, k) \in \mathrm{J} \times \mathrm{J}\) such that \(b \in \mathrm{F}_j \cap \mathrm{F}_k\). We define a section \(s' \colon \mathrm{W} \to \mathrm{E}\) by \(s'(b) = s_j(b)\) if \(b \in \mathrm{F}_j\) . We have \(\mathrm{A} \subset \mathrm{W}\) and \(s' | \mathrm{A} = s\). The section \(s'\) is continuous by prop. 4 of I, p.~19.

It remains to prove that W is a neighbourhood of A in B. For every pair \((j,k)\in\mathrm{J}\times\mathrm{J}\) , denote by \(T_{jk}\) the set of points \(b\in F_{j}\cap F_{k}\) such that \(s_{j}(b)\neq s_{k}(b)\) . The set \(T_{jk}\) is closed in \(F_{j}\cap F_{k}\) (prop. 11, b) of I, p.~34), hence in V, and the sets \(T_{jk}\) form a locally finite family of subsets of V. The union T of the family \((T_{jk})\) is therefore a closed set in V (TG, I, p.~6, prop. 4). Now W is, by definition, the complement of T in V. It is therefore a neighbourhood of A in V and consequently in B.

